\section*{Metodología}
\justifying
Al desarrollar el protocolo – o la propuesta – para el proyecto, la metodología
se constituye en el diseño del proyecto. En esta sección debe explicarse claramente
la forma en que se desarrollará el proyecto. Por lo tanto, en el protocolo, la
metodología se escribe en futuro, como una ‘promesa’ o propuesta de lo que se va
a hacer y –sobre todo – como se va a hacer. 

En resumen, la metodología debe explicar todos los pasos a seguir para desarrollar 
el proyecto. En términos muy sencillos, elaborar un protocolo –también identificado
como propuesta – requiere de una estructura para organizar cada sección del mismo.
En esta estructura, la sección conocida como metodología cumple varias funciones.
Como aspecto primordial, la metodología debe explicar el diseño del proyecto.
Primero debe esbozar la forma en que se desarrollará todo el proceso, con el
mayor número de detalles posible, pero con la posibilidad de modificarse en algunos
aspectos durante su desarrollo. Si esto sucede, la persona que desarrolla el proyecto
debe explicar claramente cuáles fueron las modificaciones y las razones de peso
que se tomaron en cuenta para variar la metodología. Como ya se aclaró líneas arriba,
todo en el protocolo son promesas planteadas por la persona que presenta
la propuesta y compromisos adquiridos.
